% Options for packages loaded elsewhere
\PassOptionsToPackage{unicode}{hyperref}
\PassOptionsToPackage{hyphens}{url}
\documentclass[
  10pt,
]{article}
\usepackage{xcolor}
\usepackage[margin=0.85in]{geometry}
\usepackage{amsmath,amssymb}
\setcounter{secnumdepth}{-\maxdimen} % remove section numbering
\usepackage{iftex}
\ifPDFTeX
  \usepackage[T1]{fontenc}
  \usepackage[utf8]{inputenc}
  \usepackage{textcomp} % provide euro and other symbols
\else % if luatex or xetex
  \usepackage{unicode-math} % this also loads fontspec
  \defaultfontfeatures{Scale=MatchLowercase}
  \defaultfontfeatures[\rmfamily]{Ligatures=TeX,Scale=1}
\fi
\usepackage{lmodern}
\ifPDFTeX\else
  % xetex/luatex font selection
\fi
% Use upquote if available, for straight quotes in verbatim environments
\IfFileExists{upquote.sty}{\usepackage{upquote}}{}
\IfFileExists{microtype.sty}{% use microtype if available
  \usepackage[]{microtype}
  \UseMicrotypeSet[protrusion]{basicmath} % disable protrusion for tt fonts
}{}
\makeatletter
\@ifundefined{KOMAClassName}{% if non-KOMA class
  \IfFileExists{parskip.sty}{%
    \usepackage{parskip}
  }{% else
    \setlength{\parindent}{0pt}
    \setlength{\parskip}{6pt plus 2pt minus 1pt}}
}{% if KOMA class
  \KOMAoptions{parskip=half}}
\makeatother
\usepackage{longtable,booktabs,array}
\usepackage{calc} % for calculating minipage widths
% Correct order of tables after \paragraph or \subparagraph
\usepackage{etoolbox}
\makeatletter
\patchcmd\longtable{\par}{\if@noskipsec\mbox{}\fi\par}{}{}
\makeatother
% Allow footnotes in longtable head/foot
\IfFileExists{footnotehyper.sty}{\usepackage{footnotehyper}}{\usepackage{footnote}}
\makesavenoteenv{longtable}
\setlength{\emergencystretch}{3em} % prevent overfull lines
\providecommand{\tightlist}{%
  \setlength{\itemsep}{0pt}\setlength{\parskip}{0pt}}
\usepackage{bookmark}
\IfFileExists{xurl.sty}{\usepackage{xurl}}{} % add URL line breaks if available
\urlstyle{same}
\hypersetup{
  pdftitle={Letter of Expectations --- Fall 2026 Second Teaching Placement},
  hidelinks,
  pdfcreator={LaTeX via pandoc}}

\title{Letter of Expectations --- Fall 2026 Second Teaching Placement}
\author{}
\date{}

\begin{document}

\begin{center}
{\LARGE\bfseries Letter of Expectations}\\[4pt]
{\large Fall 2026 Second Teaching Placement}
\end{center}

\textbf{Draft status:} Working draft for Piter Garcia, Patrick Flanagan,
and Zenon Borys to review and revise together. It is not approved,
signed, or submitted. Use the current Warner cover sheet/form if
required.

\textbf{Effective date:} Date of final signature: {[}complete after
agreement{]}

\textbf{Placement dates:} September 8--November 20, 2026
(district-confirmed). This draft records the placement context and
proposes expectations for the remainder of the placement; it does not
retroactively approve schedule changes.

\textbf{Privacy:} Keep this draft in the local EDF436 workspace and the
private Overleaf project. Do not commit the letter or upload it to the
link-accessible EDF436 or TeachingPlacement Drive folders. Do not
include student-identifying information.

\subsection{Cover Letter --- complete every applicable
field}\label{cover-letter--complete-every-applicable-field}

\begin{longtable}[]{@{}p{0.20\linewidth}p{0.36\linewidth}p{0.36\linewidth}@{}}
\toprule\noalign{}
Required field & Draft entry & Confirmation needed \\
\midrule\noalign{}
\endhead
\bottomrule\noalign{}
\endlastfoot
Candidate name & Piter Garcia & Confirm preferred name/signature form \\
Candidate phone & {[}enter{]} & Candidate to complete \\
Candidate email &
\href{mailto:pgarcia8@u.rochester.edu}{\nolinkurl{pgarcia8@u.rochester.edu}}
& Confirm best address \\
School-Based Teacher Educator & Patrick Flanagan & Confirm name and
preferred title \\
SBTE phone & {[}enter preferred school contact{]} & Confirm with
Patrick \\
SBTE email & {[}enter{]} & Confirm with Patrick \\
Inclusion SBTE & {[}N/A if no separate inclusion SBTE is assigned{]} &
Confirm with Zenon; complete contact fields if applicable \\
School & Eastridge Senior High School (EHS), East Irondequoit Central
School District & Confirm official school-name formatting \\
School address & {[}enter{]} & Confirm from school directory \\
School telephone & {[}enter{]} & Confirm from school directory \\
University-Based Teacher Educator & Zenon Borys --- Computer Science &
Confirm preferred title \\
UBTE phone & {[}enter preferred number{]} & Confirm with Zenon \\
UBTE email &
\href{mailto:zborys@warner.rochester.edu}{\nolinkurl{zborys@warner.rochester.edu}}
& Confirm best address \\
Inclusion UBTE & {[}N/A if no separate inclusion UBTE is assigned{]} &
Confirm with Zenon; complete contact fields if applicable \\
Placement & Fall 2026 second teaching placement; Computer
Science--Technology & Confirm program label \\
Signatures & Candidate, SBTE, UBTE; inclusion educators if applicable &
All required participants review the handbook and sign/date the Warner
cover sheet \\
\end{longtable}

\subsection{I. General Expectations}\label{i-general-expectations}

This nine-week placement (approximately 45 school days) is a supported
opportunity to learn the school and classroom context and to contribute
to planning, instruction, assessment, and classroom routines. The
candidate, SBTE, and UBTE will make decisions that serve students, fit
the courses and school context, reflect the candidate's readiness, and
remain workable for the SBTE. The placement is not expected to follow a
fixed progression or require the candidate to assume responsibility for
every class.

The candidate's prior Spring 2026 placement and course work provide a
foundation in computer-science lesson design, clear learning goals,
inclusive supports, assessment and feedback, classroom routines, and
reflection. For this placement, the candidate's planning work has
connected those practices to Eastridge's business and marketing courses
and Microsoft Office/365 tools. Examples prepared to date include
student-created workplace norms; Microsoft 365/Teams workflow
instruction; short marketing activities tied to the 4 Ps and marketing
functions; and accessible, visually supported lesson materials with
clear directions, pacing, and checks for understanding. These examples
describe the candidate's preparation and developing practice; the SBTE
retains authority over which materials are used and when.

All participants will follow East Irondequoit Central School District,
Eastridge, Warner School, and University policies; protect student and
family confidentiality; communicate promptly and professionally; and
seek help early when expectations or circumstances are unclear. The
Student Teaching Handbook and current course requirements govern where
this agreement does not specify a procedure.

\subsection{II. Specific Requirements and
Roles}\label{ii-specific-requirements-and-roles}

\subsubsection{A. Candidate --- Piter
Garcia}\label{a-candidate--piter-garcia}

The candidate will:

\begin{enumerate}
\def\labelenumi{\arabic{enumi}.}
\tightlist
\item
  Attend on the schedule agreed below, arrive prepared, and notify the
  SBTE and UBTE as soon as possible about an absence, delay, emergency,
  or requested schedule change using the agreed channels.
\item
  Learn the assigned courses, students' learning context, curriculum,
  classroom routines, school procedures, and approved technology systems
  through observation and active participation.
\item
  Support students individually, in small groups, and during whole-class
  learning when directed by the SBTE; use respectful, strengths-based
  communication and follow the SBTE's classroom routines.
\item
  Co-plan, co-teach, and lead lesson segments or other instructional
  responsibilities as agreed with the SBTE and UBTE. Responsibility may
  grow toward primary responsibility for one or more
  classes/preparations when the arrangement serves students and all
  parties agree it is appropriate.
\item
  Prepare and review lesson plans using the agreed school and Warner
  formats. Plans will identify learning goals, lesson sequence,
  materials/technology, accessible supports, checks for understanding,
  and likely follow-up adjustments.
\item
  Use clear, student-facing directions and reasonable pacing; plan
  supports such as visual models, vocabulary scaffolds, chunked steps,
  opportunities to think and respond, and accessible digital materials
  as appropriate to the lesson and students.
\item
  Use Microsoft Office/365, Teams, classroom files, assessment tools,
  and other school systems only within the SBTE's direction and district
  policy. The SBTE retains responsibility for official grades, student
  records, family communication, attendance, and school decisions unless
  a task is expressly delegated and permitted.
\item
  Request and use feedback; reflect on student learning and
  participation; and make agreed adjustments to instruction. Maintain a
  private, organized portfolio of plans, reflections, feedback, and
  permitted de-identified evidence for program requirements.
\item
  Maintain professional conduct and boundaries; attend relevant school
  meetings or activities when invited and appropriate; and communicate
  with students, families, colleagues, and administrators only in roles
  and channels authorized by the SBTE/school.
\item
  Not act as a substitute or assume sole responsibility for a class in
  the SBTE's absence. Follow the school's direction and the Handbook's
  restriction on substituting during clinical experiences.
\end{enumerate}

\subsubsection{B. School-Based Teacher Educator --- Patrick
Flanagan}\label{b-school-based-teacher-educator--patrick-flanagan}

The SBTE will:

\begin{enumerate}
\def\labelenumi{\arabic{enumi}.}
\tightlist
\item
  Orient the candidate to the school, assigned courses, students'
  learning context, curriculum, classroom routines, safety procedures,
  technology systems, and relevant district/school policies.
\item
  Identify meaningful opportunities for the candidate to observe,
  assist, co-plan, co-teach, support students, and gradually take on
  instructional responsibility.
\item
  Decide with the candidate which class(es), lesson(s), and
  responsibilities are appropriate at each stage; the candidate is not
  expected to lead every course or section.
\item
  Model instruction, classroom routines, assessment, feedback, and
  professional practice; clarify what the candidate may do independently
  and when approval is needed.
\item
  Review lesson plans/materials on the agreed schedule, give timely,
  actionable feedback, and identify any required changes before the
  candidate teaches or distributes materials.
\item
  Confer regularly with the candidate about student learning, upcoming
  lessons, classroom routines, feedback, and next steps; arrange
  additional conversations when needed.
\item
  Provide ongoing evidence and feedback for program evaluation, complete
  required school-based forms, and collaborate with the UBTE on
  observations, evaluation, and a thoughtful transition into and out of
  the placement.
\item
  Serve as the school liaison for student/family matters, official
  records and grading, discipline, attendance, and school decisions;
  clarify the boundaries of any delegated candidate tasks.
\end{enumerate}

\subsubsection{C. University-Based Teacher Educator --- Zenon
Borys}\label{c-university-based-teacher-educator--zenon-borys}

The UBTE will:

\begin{enumerate}
\def\labelenumi{\arabic{enumi}.}
\tightlist
\item
  Support the candidate and SBTE in finalizing and, when useful,
  revisiting this Letter of Expectations.
\item
  Clarify University/Warner requirements, the observation process,
  required forms, timelines, lesson-plan submission expectations, and
  approved submission channels.
\item
  Typically conduct three observations during the nine-week placement,
  at mutually workable times that are instructionally meaningful. The
  exact dates and observation procedures will be agreed separately.
\item
  Provide observation feedback and facilitate a post-observation
  conversation with the candidate and SBTE as soon as practical after
  each observation.
\item
  Coordinate the two CPAST administrations with the candidate and SBTE.
  The candidate completes self-assessments; the SBTE and UBTE complete
  evaluator training, independently assess the candidate, then meet to
  review evidence and reach consensus. The midpoint process is
  formative; the final process is summative. Timing and meeting format
  will be scheduled by the UBTE.
\item
  Remain available to the candidate and SBTE for questions or concerns
  about the placement, University requirements, evaluation, or support.
\end{enumerate}

\subsection{III. Schedule and Responsibility
Timeline}\label{iii-schedule-and-responsibility-timeline}

\subsubsection{A. Attendance schedule}\label{a-attendance-schedule}

\begin{longtable}[]{@{}p{0.20\linewidth}p{0.36\linewidth}p{0.36\linewidth}@{}}
\toprule\noalign{}
Schedule item & Proposed agreement & Confirm before signature \\
\midrule\noalign{}
\endhead
\bottomrule\noalign{}
\endlastfoot
Placement span & September 8--November 20, 2026; district-confirmed &
Any individually approved exception or school-closure adjustment \\
Placement days & Monday--Friday, following the live Eastridge calendar
and SBTE assignment & Which days the candidate is expected on site;
confirm any date-specific University-course conflict plan \\
Arrival & By 7:00 a.m. to prepare for the 7:30 a.m. first class;
supported by the placement planning conversation & Patrick confirms this
remains the expected arrival time \\
Daily departure & {[}agree on the expected departure time and any
different date-specific arrangement{]} & Do not treat personal calendar
travel/departure blocks as school approval \\
Class schedule & Current placement map includes Introduction to
Marketing (Period 1, 7:30--8:12 a.m.; Period 5C6A, 11:32 a.m.--12:14
p.m.) and Principles and Tools of Business (Period 4, 9:58--10:40 a.m.;
Period 7, 12:49--1:31 p.m.; Microsoft Office in Term 1) & Patrick
confirms assigned classes/sections and current bell/advisory schedule \\
Absence / urgent contact & {[}agree on same-day phone/text route and
follow-up email{]} & Patrick and Zenon confirm preferred contacts \\
Make-up time & Determined under district and University policy with the
SBTE/UBTE & Confirm how missed time is recorded or made up \\
\end{longtable}

The class list and bell times above reflect the current placement map
and school schedule reviewed for preparation. The SBTE's current
assignment and live school schedule control. The candidate will not
treat calendar planning blocks, travel buffers, or course-conflict
workarounds as school-approved attendance hours.

\subsubsection{B. Planned progression of
responsibilities}\label{b-planned-progression-of-responsibilities}

The following is a discussion framework, not a fixed week-by-week
mandate. At each stage, the candidate and SBTE will set the next step
based on students' needs, the curriculum, readiness, and the SBTE's
judgment. The UBTE will advise as needed.

\begin{longtable}[]{@{}p{0.20\linewidth}p{0.36\linewidth}p{0.36\linewidth}@{}}
\toprule\noalign{}
Stage & Candidate focus & Agreed check before moving forward \\
\midrule\noalign{}
\endhead
\bottomrule\noalign{}
\endlastfoot
Orientation and entry & Learn routines and assigned courses; observe
teaching, student engagement, assessment, and classroom expectations;
assist students and prepare bounded contributions & SBTE identifies
appropriate tasks, access boundaries, and first planning/feedback
routine \\
Supported participation & Co-plan and co-teach selected segments;
support individual/small-group/whole-class learning; use short
activities and materials approved by the SBTE & Review lesson goals,
directions, student supports, time, and checks for understanding before
use \\
Increased responsibility & Lead agreed lessons, segments, or one or more
preparations when appropriate; assess learning and reflect with the
SBTE; adjust based on evidence & SBTE and candidate agree on scope,
timing, supervision, grading/record duties, and feedback \\
Transition and exit & Return all assumed teaching and classroom
responsibilities deliberately; finish agreed
grading/documentation/follow-up; hand off materials and current
student-learning information through approved channels & Candidate,
SBTE, and UBTE agree on final responsibilities and confirm nothing
assumed by the candidate is left incomplete \\
\end{longtable}

The specific classes to observe, co-teach, and lead, and the timing of
each transition, remain {[}to be agreed with Patrick and Zenon{]}.

\subsection{IV. Lesson Plans}\label{iv-lesson-plans}

The candidate and SBTE will use {[}school lesson-plan format/location{]}
for school teaching and {[}Warner lesson-plan format when required for
an observation or course submission{]}. Before the candidate teaches or
distributes a new lesson/activity, they will review it with the SBTE at
least {[}agreed lead time{]} in advance. Routine plans will be reviewed
{[}agreed cadence{]}; the candidate and SBTE will set a regular
planning/check-in time of {[}day/time{]}.

Reviews will confirm the learning target and course alignment, lesson
sequence, anticipated student responses, accessible supports, materials
and technology, time limits/transitions, checks for understanding, and
how evidence will inform next steps. The SBTE decides whether a proposed
activity is used, adapted, or deferred. Candidate work to
date---including clear marketing-task directions, student-created
classroom-norms planning, Microsoft 365 workflow materials, and visually
supported/UDL-informed resources---may be adapted only when the SBTE
approves it for the assigned class and current lesson sequence.

\subsection{V. Meetings and
Communication}\label{v-meetings-and-communication}

\begin{enumerate}
\def\labelenumi{\arabic{enumi}.}
\tightlist
\item
  \textbf{Regular SBTE--candidate meeting:} Meet {[}weekly / other
  agreed cadence{]} at {[}day/time or planning period{]} for
  approximately {[}duration{]}. Focus: upcoming instruction, lesson-plan
  review, student-support needs, feedback, reflection, and next steps.
\item
  \textbf{Additional meetings:} Either the candidate or SBTE may request
  a meeting through {[}agreed school channel{]}. Arrange it at a
  mutually workable time; address urgent safety or student concerns
  through school procedures immediately.
\item
  \textbf{UBTE check-ins:} The candidate, SBTE, and UBTE will
  communicate as needed to coordinate the letter, observations,
  assessments, and concerns.
\item
  \textbf{Post-observation meeting:} The candidate, SBTE, and UBTE will
  meet as soon as practical after each UBTE observation to discuss
  evidence, feedback, questions, and specific next steps.
\item
  \textbf{Routine / urgent contact:} Use {[}approved email/channel{]}
  for planning and formal records and {[}approved call/text route{]} for
  same-day absence, delay, or urgent updates. Confirm which channel each
  person should use.
\item
  \textbf{Confidentiality:} Do not place identifiable student
  information in personal files, email, screenshots, presentations, or
  external services. Store and share school/student records only through
  district-approved systems and permissions.
\end{enumerate}

\subsection{VI. Observations and
Evaluation}\label{vi-observations-and-evaluation}

\begin{longtable}[]{@{}p{0.20\linewidth}p{0.36\linewidth}p{0.36\linewidth}@{}}
\toprule\noalign{}
Process & Working expectation & Confirm with UBTE/SBTE \\
\midrule\noalign{}
\endhead
\bottomrule\noalign{}
\endlastfoot
SBTE observation and feedback & Ongoing classroom feedback; pre-lesson
review when appropriate & Observation/check-in cadence and feedback
method \\
UBTE observations & Typically three during the placement & Dates,
lesson-plan lead time, and observation protocol \\
Post-observation conversation & Candidate, SBTE, and UBTE meet as soon
as practical after each observation & Scheduling method \\
CPAST & Two administrations; candidate self-assessment and evaluator
assessments by SBTE/UBTE, followed by evidence review and consensus &
Dates; midpoint and final meeting format \\
Field-placement evaluation & The syllabus describes a 100-hour
field-placement evaluation before transition to student teaching & UBTE
confirms whether and when that form applies to this second-placement
segment \\
Evidence and portfolio & Candidate maintains permitted plans,
reflections, feedback, and de-identified evidence for evaluation and the
Field Experience Resource Guide & Approved location, access, retention,
and sharing rules \\
\end{longtable}

\subsection{VII. Possible Situations}\label{vii-possible-situations}

\begin{enumerate}
\def\labelenumi{\arabic{enumi}.}
\tightlist
\item
  \textbf{SBTE absence:} The candidate contacts the designated
  administrator or school contact and follows the school's direction.
  The candidate does not substitute or take sole responsibility for the
  class.
\item
  \textbf{Candidate absence or delay:} The candidate contacts the SBTE
  and UBTE promptly through the agreed urgent channel, then follows any
  required written notice and make-up process.
\item
  \textbf{School closure, weather, or emergency:} Follow East
  Irondequoit/Eastridge announcements and procedures. The candidate
  contacts the SBTE/UBTE if the effect on attendance or instruction is
  unclear.
\item
  \textbf{Schedule or University-course conflict:} Raise the conflict
  with the SBTE and UBTE as early as possible. A changed arrival,
  departure, or missed period is effective only when the relevant school
  and University participants approve it.
\item
  \textbf{Instructional, student-support, or safety concern:} Follow
  school procedure and promptly consult the SBTE. Contact the UBTE when
  the concern needs University support or cannot be resolved through the
  regular process.
\item
  \textbf{Change in assignment or responsibility:} The candidate, SBTE,
  and UBTE will discuss the change, its effect on students and course
  requirements, and any needed revision to this letter.
\end{enumerate}

\subsection{Agreement and signatures}\label{agreement-and-signatures}

By signing the Warner cover sheet, the parties confirm that they have
reviewed the Student Teaching Handbook and agree to the responsibilities
and procedures in the final version of this letter. Complete all
applicable cover-sheet fields, replace every bracketed item with an
agreed value or an explicit N/A, and attach this agreement to the
required Warner form.

\begin{longtable}[]{@{}p{0.20\linewidth}p{0.36\linewidth}p{0.36\linewidth}@{}}
\toprule\noalign{}
Participant & Signature & Date \\
\midrule\noalign{}
\endhead
\bottomrule\noalign{}
\endlastfoot
Candidate: Piter Garcia &
\_\_\_\_\_\_\_\_\_\_\_\_\_\_\_\_\_\_\_\_\_\_\_\_\_\_\_\_\_\_ &
\_\_\_\_\_\_\_\_\_\_ \\
School-Based Teacher Educator: Patrick Flanagan &
\_\_\_\_\_\_\_\_\_\_\_\_\_\_\_\_\_\_\_\_\_\_\_\_\_\_\_\_\_\_ &
\_\_\_\_\_\_\_\_\_\_ \\
University-Based Teacher Educator: Zenon Borys &
\_\_\_\_\_\_\_\_\_\_\_\_\_\_\_\_\_\_\_\_\_\_\_\_\_\_\_\_\_\_ &
\_\_\_\_\_\_\_\_\_\_ \\
Inclusion School-Based Teacher Educator, if assigned &
\_\_\_\_\_\_\_\_\_\_\_\_\_\_\_\_\_\_\_\_\_\_\_\_\_\_\_\_\_\_ &
\_\_\_\_\_\_\_\_\_\_ \\
Inclusion University-Based Teacher Educator, if assigned &
\_\_\_\_\_\_\_\_\_\_\_\_\_\_\_\_\_\_\_\_\_\_\_\_\_\_\_\_\_\_ &
\_\_\_\_\_\_\_\_\_\_ \\
\end{longtable}

\end{document}
